\documentclass[twocolumn]{iopjournal}

\usepackage{amsmath,amssymb,amsfonts}
\usepackage{amsthm}

\fancyhfoffset[L]{0mm}
\setlength{\marginparwidth}{0pt}
\setlength{\marginparsep}{0pt}

% ==========================================
% ZENODO PREPRINT HEADERS (ACTIVE)
% ==========================================
\renewcommand{\articletype}[1]{%
  {\vspace*{-8mm}\noindent \Large \sf Preprint}\par
  \vspace*{6mm}\noindent{\scriptsize \sf{\bfseries \MakeUppercase{#1}}}}

\fancyhead[L]{}
\fancyhead[R]{}

% Theorem environments
\newtheorem{lemma}{Lemma}
\newtheorem{theorem}[lemma]{Theorem}
\newtheorem{corollary}[lemma]{Corollary}
\theoremstyle{definition}
\newtheorem{definition}[lemma]{Definition}
\newtheorem{remark}[lemma]{Remark}

\begin{document}

\articletype{Supplementary Material}

\title{The Spatial Geometry of the Lorentz Tetrad}

\author{Stephen Nelson}

\affil{Field Dynamics Research}

\ead{stephen@fielddynamics.org}

\begin{abstract}
A Lorentz tetrad is a 4-indexed object with signature
$\eta_{ab} = \mathrm{diag}(-1,+1,+1,+1)$; its index set carries the
natural action of $S_4$.  We classify its spatial realizations in
$\mathbb{R}^3$ under three conditions: common norm, balanced
closure, and $S_4$-invariance of the Gram matrix.  The centerpiece
is an orbit lemma: $S_4$-invariant $4 \times 4$ matrices form a
two-dimensional space spanned by the identity $I$ and the all-ones
matrix $J$.  The three conditions then pin both coefficients,
forcing the Gram matrix to $G = c^2 I - (c^2/3)(J - I)$.
Geometric identification gives the regular tetrahedron with
circumradius $c$, unique up to $O(3)$.  The canonical $\{-1,+1\}$
sign matrix, spatial parity, the cube $\{-1,+1\}^3$, and Kepler's
stella octangula~\cite{Kepler1619,Cromwell1997} follow as
corollaries.
\end{abstract}

% ==================================================================

\section*{The Lorentz tetrad}

A Lorentz tetrad is a 4-indexed object with index set
$\{0, 1, 2, 3\}$, where index 0 carries signature $-1$ and
indices $1, 2, 3$ carry signature $+1$ under the Lorentz metric
\begin{equation}
  \eta_{ab} = \mathrm{diag}(-1,+1,+1,+1).
  \label{eq:lorentz}
\end{equation}
The indexing is part of the object's definition: $S_4$ acts on
$\{0,1,2,3\}$, and this action carries canonically to any indexed
$4$-tuple realization of the tetrad.  The $S_4$ action is on the
index set as a set; the signature $\eta$ itself has stabilizer
$S_3 \subset S_4$.  The Lorentz signature fixes the indexed sign
pattern on the four components.

A \emph{spatial realization} of the tetrad is a $4$-tuple
$V = (v_0, v_1, v_2, v_3)$ with $v_a \in \mathbb{R}^3$, indexed by
$\{0, 1, 2, 3\}$.  We identify three conditions F1--F3 on such a
realization --- common Euclidean norm on the single diagonal
orbit, balanced closure, and $S_4$-invariance of the Gram matrix
--- and classify those realizations that satisfy all three.  The
classification shows that they are unique up to the orthogonal
group.

% ==================================================================

\section*{Realization conditions}

\begin{definition}[Spatial realization satisfying F1--F3]
\label{def:equiv}
A spatial realization $V = (v_0, v_1, v_2, v_3) \in (\mathbb{R}^3)^4$
of the Lorentz tetrad satisfies \emph{conditions F1--F3}
if:

\textup{(F1)}~\emph{Common norm}:\quad $|v_a| = c$ for all $a$,
with $c > 0$ fixed.

\textup{(F2)}~\emph{Balanced closure}:\quad
$\sum_{a=0}^{3} v_a = 0$.

\textup{(F3)}~\emph{$S_4$-invariance of the Gram matrix}:\quad the
Gram matrix $G_{ab} = v_a \cdot v_b$ satisfies
$P_\sigma G P_\sigma^T = G$ for every $\sigma \in S_4$, where
$P_\sigma$ is the permutation matrix of $\sigma$.
\end{definition}

\noindent
Condition F1 records the common Euclidean norm on the single
diagonal orbit.
Condition F2 excludes the rank-one degenerate realization
$v_0 = v_1 = v_2 = v_3$, which satisfies F1 and F3 but collapses
all four components to a single vector.  Condition F3 is the invariance
of the Gram matrix under the $S_4$ action on the index set
$\{0,1,2,3\}$; it produces the two-parameter commutant of
Lemma~\ref{lem:commutant}.

\noindent
Under F1, F2, and F3 the forced Gram matrix has rank three with
kernel $\mathrm{span}\{(1,1,1,1)\}$ (Cor.~\ref{cor:rank}).

% ==================================================================

\section*{The commutant lemma}

The algebraic heart of the classification is a single
orbit-counting fact about $S_4$-invariant $4 \times 4$ matrices.

\begin{lemma}[Commutant of the permutation representation of $S_4$]
\label{lem:commutant}
Let $\rho: S_4 \to GL_4(\mathbb{R})$ be the permutation
representation on $\mathbb{R}^4$.  A matrix $G \in M_4(\mathbb{R})$
satisfies $\rho(\sigma)\,G\,\rho(\sigma)^T = G$ for all
$\sigma \in S_4$ if and only if
\begin{equation}
  G = \alpha\,I + \beta\,J,
  \label{eq:commutant}
\end{equation}
for some scalars $\alpha, \beta \in \mathbb{R}$, where $I$ is the
$4 \times 4$ identity and $J$ is the $4 \times 4$ all-ones matrix.
\end{lemma}

\begin{proof}
The condition $\rho(\sigma)\,G\,\rho(\sigma)^T = G$ is equivalent
to $G_{\sigma(a)\sigma(b)} = G_{ab}$ for every $\sigma \in S_4$ and
every pair $(a, b) \in \{0, 1, 2, 3\}^2$.  Thus $G_{ab}$ is constant
on the orbits of the natural $S_4$ action on ordered pairs.  This
action has exactly two orbits: the diagonal pairs $\{(a,a) : a \in
\{0,1,2,3\}\}$, and the off-diagonal pairs $\{(a,b) : a \neq b\}$
(the latter is a single orbit because $S_4$ is $2$-transitive on
$\{0,1,2,3\}$).  Hence $G_{aa}$ takes a single value on the diagonal
and $G_{ab}$ takes a single value off the diagonal.  Writing these
as $\alpha + \beta$ and $\beta$ respectively gives
$G = \alpha I + \beta J$.  Conversely, any matrix of this form
commutes with every permutation matrix, so both directions hold.
\hfill$\square$
\end{proof}

\noindent
Writing $G = \alpha I + \beta J$ is equivalent to writing
$G = (\alpha + \beta)I + \beta(J - I)$: the diagonal and off-diagonal
entries of $G$ are each a single constant, namely $G_{aa} =
\alpha + \beta$ and $G_{ab} = \beta$ for $a \neq b$.

% ==================================================================

\section*{Main result}

\begin{theorem}[Forced Gram matrix]
\label{thm:forced}
Let $V$ be a spatial realization of the Lorentz tetrad
satisfying F1--F3 with common norm $c$.  Then the Gram matrix of $V$ is
\begin{equation}
  G = c^2 I - \frac{c^2}{3}(J - I),
  \label{eq:forced_gram}
\end{equation}
equivalently $G_{aa} = c^2$ and $G_{ab} = -c^2/3$ for $a \neq b$.
\end{theorem}

\begin{proof}
By F3 and Lemma~\ref{lem:commutant}, $G = \alpha I + \beta J$ for
some scalars $\alpha, \beta$.  Reading off diagonal and off-diagonal
entries,
\[
  G_{aa} = \alpha + \beta \quad (a = b), \qquad
  G_{ab} = \beta \quad (a \neq b).
\]
By F1, $G_{aa} = c^2$, so $\alpha + \beta = c^2$.

By F2,
\begin{equation}
  0 = \left|\sum_{a} v_a\right|^2
    = \sum_{a,b} G_{ab}
    = 4 c^2 + 12 \beta,
  \label{eq:closure_beta}
\end{equation}
where the last equality uses four diagonal entries of value $c^2$
and twelve off-diagonal entries of value $\beta$.  Hence
$\beta = -c^2/3$, and $\alpha = c^2 - \beta = 4c^2/3$.  Substituting,
$G = (4c^2/3)I - (c^2/3)J = c^2 I - (c^2/3)(J - I)$.
\hfill$\square$
\end{proof}

% ==================================================================

\section*{Corollaries}

\begin{corollary}[Rank and spectrum]
\label{cor:rank}
The forced Gram matrix~(\ref{eq:forced_gram}) has spectrum
$\{0, \tfrac{4c^2}{3}, \tfrac{4c^2}{3}, \tfrac{4c^2}{3}\}$ and rank
three.  The null eigenvector is $(1,1,1,1)$; the eigenvalue $4c^2/3$
has multiplicity three on its orthogonal complement.
\end{corollary}

\begin{proof}
$J$ has spectrum $\{4, 0, 0, 0\}$ with eigenvector $(1,1,1,1)$ for
$\lambda = 4$.  From~(\ref{eq:forced_gram}),
$G = (4c^2/3)I - (c^2/3)J$, so the eigenvalues of $G$ are
$4c^2/3 - (c^2/3)\lambda_J$ for each eigenvalue $\lambda_J$ of $J$:
$4c^2/3 - (c^2/3)(4) = 0$ on $(1,1,1,1)$, and $4c^2/3$ on the
orthogonal complement.  Rank three follows.
\hfill$\square$
\end{proof}

\begin{corollary}[Regular tetrahedron]
\label{cor:tet}
A spatial realization $V$ satisfying F1--F3 is the vertex set of a regular
tetrahedron with circumradius $c$ and edge length $c\sqrt{8/3}$,
unique up to $O(3)$.
\end{corollary}

\begin{proof}
All six pairwise squared distances are equal:
\begin{equation}
  |v_a - v_b|^2 = G_{aa} + G_{bb} - 2 G_{ab}
                = 2c^2 + \frac{2c^2}{3} = \frac{8c^2}{3}.
  \label{eq:edge}
\end{equation}
Four equidistant points in $\mathbb{R}^3$ are the vertices of a
regular tetrahedron~\cite[\S3.7]{Coxeter1973}.  By
Corollary~\ref{cor:rank} the Gram matrix has rank three, so any two
realizations with the Gram
matrix~(\ref{eq:forced_gram}) differ by an element of
$O(3)$~\cite[Thm.~7.3.11]{HornJohnson2013}.
\hfill$\square$
\end{proof}

\begin{corollary}[Canonical $\{-1,+1\}$ representative]
\label{cor:canonical}
With $c = \sqrt{3}$, the columns of
\begin{equation}
  M = \begin{pmatrix}
    +1 & +1 & -1 & -1 \\
    +1 & -1 & +1 & -1 \\
    +1 & -1 & -1 & +1
  \end{pmatrix}
  \label{eq:M}
\end{equation}
are a spatial realization satisfying F1--F3: $|v_a|^2 = 3$ and $v_a \cdot v_b = -1$
for $a \neq b$, matching~(\ref{eq:forced_gram}) at $c^2 = 3$.  By
Corollary~\ref{cor:tet}, every realization satisfying F1--F3 with $c = \sqrt{3}$
is $O(3)$-equivalent to $M$.
\end{corollary}

\begin{proof}
Each column has entries in $\{-1,+1\}$, so $|v_a|^2 = 3$.  Each pair
of columns agrees in exactly one of three coordinates (by inspection
of~(\ref{eq:M})), giving $v_a \cdot v_b = 1 - 1 - 1 = -1$.  The
Gram matrix is therefore the $c^2 = 3$ case
of~(\ref{eq:forced_gram}).
\hfill$\square$
\end{proof}

% ==================================================================

\section*{Parity completion and the stella octangula}

\begin{corollary}[Cube as parity completion]
\label{cor:cube}
Let $V = \{v_0, v_1, v_2, v_3\}$ be the canonical representative of
Corollary~\ref{cor:canonical}, and let $P(V) = \{-v_0, -v_1, -v_2,
-v_3\}$.  Then $V \cup P(V)$ is the vertex set of the cube
$\{-1,+1\}^3$.
\end{corollary}

\begin{proof}
The columns of $M$ have sign patterns
$(+,+,+), (+,-,-), (-,+,-), (-,-,+)$.  Their negatives have sign
patterns $(-,-,-), (-,+,+), (+,-,+), (+,+,-)$.  These eight triples
are pairwise distinct and exhaust the $2^3 = 8$ elements of
$\{-1,+1\}^3$.
\hfill$\square$
\end{proof}

\begin{theorem}[Stella octangula]
\label{thm:stella}
The union $V \cup P(V)$ of Corollary~\ref{cor:cube}, viewed as the
compound of the two regular tetrahedra with vertex sets $V$ and
$P(V)$, is Kepler's stella
octangula~\cite{Kepler1619,Cromwell1997}.
\end{theorem}

\begin{proof}
Both $V$ and $P(V)$ have Gram matrix~(\ref{eq:forced_gram})
(negation preserves inner products), so each is a regular
tetrahedron by Corollary~\ref{cor:tet}.  They are antipodal and
share a center at the origin.  By Corollary~\ref{cor:cube} they are
inscribed in $\{-1,+1\}^3$.  The compound of two dual regular
tetrahedra inscribed in a cube is Kepler's stella
octangula~\cite{Kepler1619,Cromwell1997}.
\hfill$\square$
\end{proof}

% ==================================================================

\section*{Remarks}

\begin{remark}[Inevitability from two parameters]
\label{rem:two_params}
The algebraic heart of the classification is that
$S_4$-equivariance (F3) reduces the Gram matrix from a generic
symmetric object (10 parameters) to an element of a two-dimensional
commutant spanned by $I$ and $J$ (Lemma~\ref{lem:commutant}).
Common norm (F1) fixes one parameter; balanced closure (F2) fixes
the other.  The regular tetrahedron is the unique configuration
consistent with these three conditions.  Everything downstream ---
the edge length, the rank of $G$, the sign representative, the
cube, the stella octangula --- is bookkeeping.
\end{remark}

\begin{remark}[What each condition does]
\label{rem:roles}
F1 and F2 are constraints that pin scalar coefficients in the
commutant; F3 is the structural condition that produces the
commutant in the first place.  Without F3, the Gram matrix is a
generic $4 \times 4$ symmetric matrix (10 parameters) and common
norm plus closure leave eight degrees of freedom ---
insufficient to force any geometric configuration.  F3 is the
load-bearing condition for uniqueness.
\end{remark}

\begin{remark}[Invariance forbids slot-dependent magnitudes]
\label{rem:perturbation}
$S_4$-invariance of the Gram matrix forbids slot-specific diagonal
data in a spatial realization.  Consider, for contrast, a
hypothetical Gram matrix with diagonal $(1.03, 1, 1, 1)$: the
transposition $0 \leftrightarrow 1$ sends this to the matrix with
diagonal $(1, 1.03, 1, 1)$, which is not the same matrix.  The
perturbed assignment is therefore a different, symmetry-broken
Gram matrix, not an alternative realization of the same tetrad.
Under F3, the diagonal entries of the Gram matrix lie in a single
orbit of the $S_4$ action, forcing all four to take the common
value $G_{aa} = c^2$.  Invariance rules out any slot-dependent
perturbation without appeal to the specific numerical value of
$c$; the normalization $c$ itself is fixed separately by F1.
\end{remark}

\begin{remark}[Conditions are structural, not arbitrary]
\label{rem:structural}
General relativity fixes the tetrad's indexed Lorentz structure;
F1--F3 are the corresponding spatial realization conditions.  F1
fixes the common diagonal norm, F2 fixes the balanced closure
relation, and F3 fixes the $S_4$-invariant Gram structure.
Theorem~\ref{thm:forced} then forces the regular tetrahedron,
with parity completion yielding the cube and Kepler's stella
octangula~\cite{Kepler1619,Cromwell1997}.
\end{remark}

% ==================================================================

\begin{thebibliography}{9}

\bibitem{Kepler1619}
J.~Kepler,
\textit{Harmonices Mundi}
(Linz, 1619).

\bibitem{Coxeter1973}
H.~S.~M.~Coxeter,
\textit{Regular Polytopes},
3rd ed.\ (Dover, New York, 1973).

\bibitem{Cromwell1997}
P.~R.~Cromwell,
\textit{Polyhedra}
(Cambridge University Press, Cambridge, 1997).

\bibitem{HornJohnson2013}
R.~A.~Horn and C.~R.~Johnson,
\textit{Matrix Analysis},
2nd ed.\ (Cambridge University Press, Cambridge, 2013).

\end{thebibliography}

\end{document}
